\begin{table}[htbp]
\caption{Complete \textit{Pleurotus} mitogenomes, all annotated with the pipeline used here rather than with their deposited features.}\label{tab:comparative}
\footnotesize
\begin{tabular}{@{}llrrrrrr@{}}
\toprule
Species / strain & Accession & Length (bp) & GC\% & tRNA & ORFs & gpI hits & Intron bp \\
\midrule
P. citrinopileatus & NC\_036998.1 & 60,694 & 29.0 & 24 & 51 & 2 & 933 \\
P. ostreatus PC9.15 & CM057219.1 & 65,277 & 26.2 & 25 & 47 & 3 & 523 \\
P. ostreatus N001 & PX724300.1 & 65,303 & 26.2 & 25 & 47 & 4 & 636 \\
P. ostreatus & OR030114.1 & 65,311 & 26.2 & 25 & 46 & 4 & 636 \\
P. pulmonarius & NC\_061177.1 & 70,671 & 25.3 & 25 & 49 & 6 & 1,367 \\
P. ostreatus BOM\_ss14 & CM148778.1 & 71,947 & 26.6 & 25 & 51 & 5 & 1,295 \\
P. ostreatus BOM\_ss5 (this study) & CM148777.1 & 71,949 & 26.6 & 25 & 51 & 4 & 1,123 \\
P. ostreatus F515 & PX724301.1 & 72,132 & 26.7 & 25 & 54 & 7 & 1,783 \\
P. cornucopiae & NC\_038091.1 & 72,134 & 26.7 & 25 & 53 & 7 & 1,783 \\
P. ostreatus N009 & PX724302.1 & 73,085 & 26.7 & 25 & 55 & 8 & 1,912 \\
P. platypus & NC\_036999.1 & 73,807 & 26.8 & 25 & 55 & 8 & 1,912 \\
P. giganteus & NC\_062374.1 & 102,950 & 25.0 & 25 & 86 & 12 & 2,994 \\
\bottomrule
\end{tabular}
\begin{tablenotes}[flushleft]\footnotesize\item gpI hits are Rfam group I covariance-model matches; the span they cover is a lower bound, since the models capture the conserved core rather than the full intron.\end{tablenotes}
\end{table}
